% page 126 of the facsimile (image p-125 of the scan)
% block 154.9 x 237.7 mm ; left margin 8.0 mm
\pgc{-12.700mm}{0.000mm}{
\l{\cel{8}116.}
\l{}
\l{\sou{divergar}. (\sou{netrans.}) (\sou{aludante elementi plu o min inter-apuda}}
\l{\cel{3}an\cel{1}\sou{la departo-punto}). Inter-eskartar sempre plu multe,}
\l{\cel{3}segun ke li plulongeskas. - DEFIS.}
\l{}
\l{\sou{diver}sa. I. (\sou{aludante kozi quin onu kompar}as). De qua singla}
\l{\cel{3}havas karaktero diferanta de olta di la ceteri. - II.}
\l{\cel{3}Plura (personi o kozi) di qui la naturo, la karaktero}
\l{\cel{3}esas interdiferanta. - DEFIRS.}
\l{}
\l{\sou{divertikulo}.\cel{1}\sou{(anat}.)Apendico kava, analoga a sako-strado,}
\l{\cel{3}qua debushas en kavajo naturala di la homo-korpo.}
\l{}
\l{diversionar. (\sou{netrans}.)I. (\sou{pri operaco militala}).Deturnar}
\l{\cel{3}la enemiko de la punto quan lu atakas, per koaktar lu a}
\l{\cel{3}su-defenso en punto altra. -II.\cel{2}Ago qua deturnas la spi-}
\l{\cel{3}rito, la kordio, de to quo fatigas lu, suciigas lu, ad}
\l{\cel{3}objekto altra. - DEFIRS.}
\l{}
\l{\sou{dividar}. (\sou{trans.}) I. (\sou{matem.})\cel{2}Donite la produto de du fak-}
\l{\cel{3}tori ed un ek ici, determinar la faktoro nekonocata.-}
\l{\cel{3}II. Igar, de toto, plura parti. - DEFIRS.}
\l{}
\l{\sou{divinar}. (\sou{trans}.)I. Savar, per moyeni supernatura, to quo}
\l{\cel{3}celesas en la tempo pasinta, prezenta, futura. - II.}
\l{\cel{3}Deskovrar, per konjekti, supozi, interpreti, to quon onu}
\l{\cel{3}ne savas. - EFIS.}
\l{}
\l{\sou{diviziono}. I. En armeo, grupo de armeani qua kompozesas ek}
\l{\cel{3}adminime du brigadi. II. En navaro, parto de eskadro qua}
\l{\cel{3}kompozesas ek adminime tri milito-navi komandata da un}
\l{\cel{3}chefo. - DEFIRS.}
\l{}
\l{\sou{divorcar}. (\sou{netrans., de)}\cel{2}Dissolvar legale la mariajo per}
\l{\cel{3}qua onu esas unionita a la spozo. - EFIS.}
\l{}
\l{\sou{divulgar}. (\sou{trans}.)Dissavigar da multa personi. - DEFIS.}
\l{}
\l{\sou{dizastro}. Desfortuno tre granda e subita, qua ruinas multa}
\l{\cel{3}homi o tota socio, exemple : grava vinkeso, granda in-}
\l{\cel{3}cendio, naufrajo. - EFIS.}
\l{}
\l{\sou{\textquotedbl{}djin\textquotedbl{}}.- (\sou{che la Arabi}) . Demono.}
\l{}
\l{\sou{\textquotedbl{}do\textquotedbl{}}. (\sou{Muziko}). Noto unesma di la gamo ordinara, sen aci-}
\l{\cel{3}dento.}
\l{}
\l{do.\cel{2}Konjunciono, uzata por kunduktar la konsequanto, la}
\l{\cel{3}konkluzo di lo preiranta. - FI.}
\l{}
\l{\sou{doario}. IJiuro olima). Parto de la havaji donita da la spo-}
\l{\cel{3}zulo a sua spozino, quan elu posedeskis se elu divenis}
\l{\cel{3}vidva, e qua transpasis ad ilua filii. - II. Revenuo}
\l{\cel{3}atribuata a rejino, a princino salika. - FI.}
}
